\section{A global Killing obstruction invisible on every positive angular neighborhood}
\label{sec:r13-punctured-three-fan-obstruction}

The following is a supplied finite coefficient example.  It has normal
fan dimension three, positive balancing, connected adjacency through
facets, bounded homogeneous finite-PA primitives on every positive
angular neighborhood, and zero branchwise Saint--Venant defect.  Its
complete global gauge equation is nevertheless inconsistent.  No
realization by an original analytic zero or by an SNC residue model
is asserted.

\subsection{Five rational hyperplanes and one exact stress}

In \(N=\mathbf R^4\), with coordinates \(x_0,x_1,x_2,x_3\), put
\[
 P_i=\ker(x_0+i x_1+i^2x_2+i^3x_3),\qquad 0\leq i\leq4,
 \qquad S_{ij}=P_i\cap P_j.
 \tag{PK1}
\]
Every four defining rows are independent, by the Vandermonde
formula.  Thus every \(P_i\) has dimension three, every
\(S_{ij}\) has dimension two, every triple intersection has
dimension one, and no nonzero point belongs to four of the planes.
Give their union the common fan induced by the five central
hyperplanes, followed by a finite rational simplicial refinement.
Denote this fan by \(\Sigma\).  Retain all five source labels.
The complete plane \(P_i\), with its induced lattice and weight
one, is balanced; therefore their sum is a positive balanced
three-cycle.  Its maximal cones are connected through facets,
since the complete fan in each plane is facet-connected and
every pair of planes meets in a two-plane.

Use the basis
\[
 p_{i1}=(-i,1,0,0),\quad
 p_{i2}=(-i^2,0,1,0),\quad
 p_{i3}=(-i^3,0,0,1)
 \tag{PK2}
\]
of \(P_i\).  For \(\Omega_i\in\bigwedge^2 P_i^*\), let
\(\omega_i\in\mathbf R^3\) be its coefficients on the pairs
\((p_{i1},p_{i2}),(p_{i1},p_{i3}),(p_{i2},p_{i3})\).
A common ordered basis of \(S_{ij}\), symmetric in \(i,j\), is
\[
 \begin{split}
 a_{ij}&=(ij,-i-j,1,0),\\
 b_{ij}&=(ij(i+j),-i^2-ij-j^2,0,1).
 \end{split}
 \tag{PK3}
\]
Consequently
\[
 \Omega_i(a_{ij},b_{ij})
       =\ell_{ij}\cdot\omega_i,
 \qquad
 \ell_{ij}=(i^2+ij+j^2,-i-j,1).
 \tag{PK4}
\]
Define the edge restriction map
\[
 D:(\mathbf R^3)^5\longrightarrow\mathbf R^{10},
 \qquad(D\omega)_{ij}
       =\ell_{ij}\cdot(\omega_i-\omega_j),\quad i<j.
 \tag{PK5}
\]

\begin{lemma}[The complete global edge equation]
\label{lem:r13-five-plane-edge-stress}
The map \(D\) has rank nine.  In the edge order
\(01,02,03,04,12,13,14,23,24,34\), its cokernel is detected by
\[
 s=(1,-3,3,-1,6,-8,3,6,-3,1).
 \tag{PK6}
\]
Thus \(D\omega=c\) has a solution if and only if
\[
                    \sum_{i<j}s_{ij}c_{ij}=0.
 \tag{PK7}
\]
For any subset of at most three plane labels, the corresponding
edge restriction map is surjective.
\end{lemma}

\begin{proof}
Set \(r_i=(i^3,-i^2,i)\).  Then
\(\ell_{ij}=(r_i-r_j)/(i-j)\).  The kernel of \(D\) consists
of the families satisfying
\((r_i-r_j)\cdot(\omega_i-\omega_j)=0\) on all edges.
The first four \(r_i\) are affinely independent.  There is a
unique affine map \(r\mapsto Ar+b\) taking their values to
the first four \(\omega_i\).  Put \(v_i=r_i-r_0\) for
\(1\leq i\leq3\).  The six edge equations among these four
vertices say
\(v_i^TAv_i=0\) and
\((v_i-v_j)^TA(v_i-v_j)=0\).  Their symmetric parts imply
\(A+A^T=0\), since the \(v_i\) are a basis.  The four equations
joining the fifth vertex to the first four then give
\[
       (r_4-r_i)\cdot(\omega_4-Ar_4-b)=0,
                          \qquad 0\leq i\leq3.
\]
Those differences span \(\mathbf R^3\), so
\(\omega_4=Ar_4+b\).  Conversely every such alternating
\(A\) and arbitrary \(b\) is in the kernel.  Its dimension
is six, and \(\operatorname{rank}D=15-6=9\).

Let \(\lambda=(1,-4,6,-4,1)\).  The fourth-difference identity
gives \(\sum_i\lambda_i=0\) and
\(\sum_i\lambda_i r_i=0\).  Therefore
\[
 \sum_{i<j}\lambda_i\lambda_j
       (r_i-r_j)\cdot(\omega_i-\omega_j)=0.
\]
The coefficients
\(\lambda_i\lambda_j(i-j)/4\) are exactly PK6.  They give
one nonzero left kernel vector; the rank calculation makes
this the whole left kernel, proving PK7.

For at most two labels surjectivity is immediate.  For three
labels \(i<j<k\), the two rows \(\ell_{ij},\ell_{ik}\)
are independent: their last coordinates are one and their
second coordinates are distinct.  Choose \(\omega_i\) with
these two row values \(c_{ij},c_{ik}\).  Choose
\(\omega_j\) with row values zero on \(\ell_{ij}\) and
\(c_{jk}\) on \(\ell_{jk}\), and take \(\omega_k=0\).
These choices solve the three equations.
\end{proof}

\subsection{Finite-PA realization of arbitrary edge data}

Choose, for every ordered pair \(i\ne j\), an alternating form
\(\theta_{ij}\in\bigwedge^2S_{ij}^*\).  Define
\[
 c_{ij}=(\theta_{ij}-\theta_{ji})(a_{ij},b_{ij}),\qquad i<j.
 \tag{PK8}
\]
There is no relation among these ten numbers: one may specify
\(\theta_{ij}\) for \(i<j\) and take \(\theta_{ji}=0\).

\begin{lemma}[Prescribed tangential rotation traces]
\label{lem:r13-plane-trace-realization}
For each \(i\), there is a continuous homogeneous finite-PA
covector field \(b_i:P_i\to P_i^*\), with \(b_i(0)=0\), such that
\[
             b_i(x)|_{S_{ij}}=\theta_{ij}(x,\cdot)
                    \quad(x\in S_{ij},\ j\ne i).
 \tag{PK9}
\]
The fields can use one common finite rational fan refinement.
They have linear closed-cell extensions and bounded first
derivatives.  Their distributional strains
\(Z_i=e(b_i)\) are bounded piecewise constant tensors.
\end{lemma}

\begin{proof}
Take a common rational simplicial refinement of the arrangement
fan.  At a ray with a chosen generator \(v\in P_i\), impose
on its covector value the restrictions
\(\theta_{ij}(v,\cdot)\) for all \(j\) such that
\(v\in S_{ij}\).  There are at most two such labels.  With two,
the two planes \(S_{ij},S_{ik}\) intersect in the line
\(\mathbf Rv\), and both prescribed covectors are zero on
that line.  They therefore extend uniquely to a covector on
\(P_i=S_{ij}+S_{ik}\).  With one or no label, choose any
extension.  Use that one value for every cone containing the ray.

On each simplicial cone extend the assigned ray values linearly.
These extensions agree on shared faces and hence give a
continuous homogeneous finite-PA field.  On a cone contained
in \(S_{ij}\), its ray values have exactly the tangential
restriction in PK9.  Linearity proves PK9 on the whole cone.
Every coefficient has a linear extension across its closed cell.
Continuity removes the coefficient-wall jump measure from its
first distributional derivative, giving the assertion about
\(e(b_i)\).
\end{proof}

All tangential strains on \(S_{ij}\) are zero, since PK9 is a
Killing field there.  Thus these prescribed face strains agree
between the source planes.  Moreover every branchwise
Saint--Venant operator vanishes distributionally, including all
coefficient walls and the origin, because \(Z_i=e(b_i)\) on
the entire vector space \(P_i\).

\subsection{Every noncentral local gauge exists; the global one need not}

\begin{theorem}[Local normalized solvability does not force JF8]
\label{thm:r13-positive-angular-not-global-gauge}
For the strains \(Z_i\) constructed above, every nonzero point
of \(|\Sigma|\) has a conical angular neighborhood on which
there is a continuous finite-PA, fully incident-compatible
primitive, with central value zero and a linear radial bound.
A finite collection of these neighborhoods covers the link.

A global continuous finite-PS, fully incident-compatible
primitive of the same strains and central value exists if and
only if PK7 holds.  In particular, the choice
\[
                 c_{01}=1,\qquad c_{ij}=0\text{ otherwise}
 \tag{PK10}
\]
has every stated local primitive but has no global primitive.
Its complete JF7 condition holds on every face, whereas JF8 fails.
\end{theorem}

\begin{proof}
Fix a nonzero point \(x\) and let \(I(x)\) be its plane labels.
There are at most three.  A sufficiently small conical angular
neighborhood meets no other source plane.  The last assertion
of the edge-map lemma supplies forms \(\Omega_i\),
\(i\in I(x)\), solving all the edge equations among those
labels.  On this neighborhood use
\[
                       b_i^{\mathrm{loc}}(n)
                                  =b_i(n)-\Omega_i(n,\cdot).
 \tag{PK11}
\]
They have the same strains, vanish at the origin, and have
bounded first derivatives and an \(O(|n|)\) bound.
The edge equations and PK9 say that their tangential
restrictions agree on every \(S_{ij}\).

For completeness, this pairwise statement gives the complete
raw compatibility in this particular arrangement.  For two
planes it is the usual extension criterion on their intersection.
For three, extend the first two compatible covectors to one
covector on \(N=P_i+P_j\).  Its difference from the third
covector vanishes on \(S_{ik}\) and \(S_{jk}\).  These are
distinct two-planes spanning \(P_k\), so the difference is
zero.  Repeated sector copies in one plane use exactly the
same full covector.  Thus all simultaneous incident relations
are satisfied.  Compactness of a finite polyhedral link gives
a finite angular cover; its coefficient and angular subdivisions
can be chosen in common.  This proves the local assertion.

If PK7 holds, choose a global solution of PK5 and use PK11
on all five planes.  The same argument gives full compatibility
away from the origin.  At the origin all covectors are zero.
The JF extension theorem supplies an ambient finite-PA extension.

Conversely suppose a global compatible primitive exists.  On
each whole plane its sector restrictions agree as full covectors
at every internal face.  They therefore define a continuous
finite-PS covector \(g_i\) on \(P_i\).  Its full distributional
strain is \(Z_i\), since no coefficient jump measure is present.
Hence \(b_i-g_i\) has zero strain on the entire plane and is
an affine Killing field.  The common value at zero removes
its translation, leaving one constant form
\(\Omega_i\in\bigwedge^2P_i^*\), independent of all sector
subdivisions.  Compatibility along every \(S_{ij}\) gives
exactly \(D\omega=c\).  The edge-map lemma now forces PK7.
For PK10 its left side is one, a contradiction.

Finally, on a common fan subdivision all the supplied \(b_i\)
are homogeneous linear on each cone.  Every complete raw
incident defect is consequently a fixed normal-linear map on
each closed face.  This is JF7, including its full distributional
Hessian formulation and every internal coefficient wall.  The
inconsistency just proved is precisely the remaining JF8
condition.
\end{proof}

\subsection{The fixed-fan real current is also closed}

Identify a symmetric strain with its real superform of bidegree
\((1,1)\).  Let \(\iota_i:P_i\hookrightarrow N\) be the
closed linear embedding, with the lattice calibration and weight
one already specified.  Define the ambient current
\[
                   \mathcal T=\sum_{i=0}^4\iota_{i*}[Z_i],
                      \qquad D=d'd''.
 \tag{PK12}
\]
Here \([Z_i]\) is the current of the locally bounded intrinsic
form on the entire plane \(P_i\).  Since \(\dim P_i=3\)
and \(\dim N=4\), the ambient current has bidegree \((2,2)\).

\begin{proposition}[Closure against all fixed-fan real smooth tests]
\label{prop:r13-five-plane-coarse-current-closed}
The current \(\mathcal T\) satisfies \(D\mathcal T=0\)
against every compactly supported smooth real superform on
\(N\) of the complementary bidegree.  This remains true for
PK10, although its global compatible strain primitive does not
exist.
\end{proposition}

\begin{proof}
On each whole plane, write \(b_i=\sum_j b_{ij}\,dx_j\)
in linear coordinates and put
\[
 \beta_i=\tfrac12\sum_j b_{ij}\,d''x_j,
 \qquad
 \zeta_i=-\tfrac12\sum_j b_{ij}\,d'x_j.
\]
With the usual superform sign convention the full distributional
identity is
\(Z_i=d'\beta_i+d''\zeta_i\).  Its coefficients are those
of the symmetric derivative of \(b_i\); continuity of \(b_i\)
removes first-derivative jump measures.  The bicomplex identities
give \(D[Z_i]=0\), with every further derivative taken in
distributions.  The full plane has no boundary.

For a compactly supported smooth test form on \(N\), its
restriction to \(P_i\) is compactly supported and restriction
commutes with both derivatives.  The definition of current
derivatives therefore gives
\(D\iota_{i*}[Z_i]=\iota_{i*}D[Z_i]=0\).  Summing proves
the proposition.  All coefficient walls are included in this
calculation; no pointwise differentiation across them is used.
\end{proof}

This statement concerns real smooth tests on this fixed fan.
It does not assert an actual analytic current identity on every
original refinement, an analytic zero relation, or equality
against GK delta tests.  It shows that fixed-fan real current
closure plus all the positive angular primitives in the theorem
still does not imply JF8.

\paragraph{Ordinary topology does not remove this coefficient obstruction.}
The link \(X=|\Sigma|\cap S^3\) has
\(H_1(X;\mathbf R)=0\).  Indeed it is the union of five
2-spheres.  When a new sphere is added, its intersection with
the preceding union is a nonempty connected union of great
circles: any two of those circles meet.  On a common finite
triangulation, the Mayer--Vietoris sequence and induction give
zero first homology.  Nevertheless the normalized Killing
coefficient sheaf of KD has a nonzero descent class for PK10.
The obstruction belongs to that varying coefficient sheaf,
not to the constant-coefficient first homology of the link.

\paragraph{Consequence for the analytic program.}
A theorem supplying compatible normalized angular primitives
on one fixed actual tuple proves the local part of the joint
criterion.  To infer its global gauge one must additionally
show that the actual data annihilate the finite Killing stresses,
or prove a geometric vanishing theorem for the actual descent
class.  Positivity, balancing, facet connectivity, bounded radial
regularity, branchwise Saint--Venant compatibility, and local
solvability on every positive ray do not alone imply that
vanishing.  PK10 is a finite coefficient diagnostic; it is not
an actual analytic relation or a counterexample to A1.
